% concatenated from PCINJDIM.tex
\documentclass[11pt,a4paper]{amsart}
\usepackage{rotating}
\usepackage[all]{xy}
\usepackage[utf8]{inputenc}
\usepackage{amscd,amssymb,amsopn,amsmath,amsthm,graphics,amsfonts,enumerate,verbatim,calc}

% Page geometry
\textwidth=15cm
\textheight=21cm
\topmargin=0.00cm
\oddsidemargin=0.00cm
\evensidemargin=0.00cm
\headheight=14.4pt
\headsep=1cm
\numberwithin{equation}{section}

% Hyphenation and spacing
\hyphenation{semi-stable}
\emergencystretch=10pt

% Footnoting
\def\thefootnote{*}
\renewcommand{\subjclassname}{%
	\textup{2020} Mathematics Subject Classification}

% Theorem-like environments
\newtheorem{theorem}{Theorem}[section]
\newtheorem{lemma}[theorem]{Lemma}
\newtheorem{proposition}[theorem]{Proposition}
\newtheorem{corollary}[theorem]{Corollary}
\newtheorem{notation}[theorem]{Notation}
\newtheorem{problem}[theorem]{Problem}
\newtheorem{observation}[theorem]{Observation}
\newcommand{\up}[1]{{{}^{#1}\!}}
\theoremstyle{definition}
\newtheorem{claim}[theorem]{Claim}
\newtheorem{definition}[theorem]{Definition}
\theoremstyle{remark}
\newtheorem{remark}[theorem]{Remark}
\newtheorem{fact}[theorem]{Fact}
\newtheorem{example}[theorem]{Example}
\newtheorem{question}[theorem]{Question}\newtheorem{discussion}[theorem]{Discussion}
\newtheorem{conjecture}[theorem]{Conjecture}
\newtheorem{answer}[theorem]{Answer}
\newtheorem{acknowledgement}{Acknowledgement}
\renewcommand{\theacknowledgement}{}

% Custom commands
\newcommand{\wtd}{\widetilde}
\newcommand{\Ass}{\operatorname{Ass}}
\newcommand{\Syz}{\operatorname{Syz}}
\newcommand{\grade}{\operatorname{grade}}
\newcommand{\Min}{\operatorname{Min}}
\newcommand{\red}{\operatorname{red}}
\newcommand{\Spec}{\operatorname{Spec}}
\newcommand{\spec}{\operatorname{spec}}
\newcommand{\h}{\operatorname{h}}
\newcommand{\Gpd}{\operatorname{Gpd}}
\newcommand{\Tr}{\operatorname{Tr}}
\newcommand{\rad}{\operatorname{rad}}
\newcommand{\cd}{\operatorname{cd}}
\newcommand{\Char}{\operatorname{Char}}
\newcommand{\End}{\operatorname{End}}
\newcommand{\Ht}{\operatorname{ht}}
\newcommand{\id}{\operatorname{id}}
\newcommand{\Cl}{\operatorname{Cl}}
\newcommand{\Rfd}{\operatorname{Rfd}}
\newcommand{\Gid}{\operatorname{Gid}}
\newcommand{\pd}{\operatorname{pd}}
\newcommand{\Gdim}{\operatorname{Gdim}}
\newcommand{\gdim}{\operatorname{gd}}
\newcommand{\Gfd}{\operatorname{Gfd}}
\newcommand{\E}{\operatorname{E}}
\newcommand{\Var}{\operatorname{Var}}
\newcommand{\injdim}{\operatorname{injdim}}
\newcommand{\Ext}{\operatorname{Ext}}
\newcommand{\Supp}{\operatorname{Supp}}
\newcommand{\Cone}{\operatorname{Cone}}
\newcommand{\Tor}{\operatorname{Tor}}
\newcommand{\Hom}{\operatorname{Hom}}
\newcommand{\RHom}{\operatorname{RHom}}
\newcommand{\Zd}{\operatorname{Zd}}
\newcommand{\Ann}{\operatorname{Ann}}
\newcommand{\Rad}{\operatorname{Rad}}
\newcommand{\wdim}{\operatorname{w.dim}}
\newcommand{\edepth}{\operatorname{Edepth}}\newcommand{\egrade}{\operatorname{Egrade}}
\newcommand{\F}{\operatorname{F}}
\newcommand{\depth}{\operatorname{depth}}
\newcommand{\pdepth}{\operatorname{p.depth}}
\newcommand{\Coass}{\operatorname{Coass}}
\newcommand{\Ker}{\operatorname{Ker}}
\newcommand{\Coker}{\operatorname{Coker}}
\newcommand{\im}{\operatorname{im}}
\newcommand{\Cosupp}{\operatorname{Cosupp}}
\newcommand{\Max}{\operatorname{Max}}
\newcommand{\ux}{\underline{x}}
\newcommand{\bac}{\backslash}
\newcommand{\al}{\alpha}
\newcommand{\BN}{\mathbb{N}}
\newcommand{\BZ}{\mathbb{Z}}
\newcommand{\vpl}{\operatornamewithlimits{\varprojlim}}
\newcommand{\vli}{\operatornamewithlimits{\varliminf}}
\newcommand{\vil}{\operatornamewithlimits{\varinjlim}}
\newcommand{\bcu}{\bigcup\limits}
\newcommand{\bca}{\bigcap\limits}
\newcommand{\botm}{\bigotimes\limits}
\newcommand{\bopm}{\bigoplus\limits}
\newcommand{\uy}{\underline{y}}
\newcommand{\lo}{\longrightarrow}
\newcommand{\fm}{\mathfrak{m}}
\newcommand{\fp}{\mathfrak{p}}
\newcommand{\fq}{\mathfrak{q}}
\newcommand{\fa}{\mathfrak{a}}
\newcommand{\fb}{\mathfrak{b}}
\newcommand{\fc}{\mathfrak{c}}
\newcommand{\fx}{\underline{x}}
\newcommand{\fn}{\mathfrak{n}}
\newcommand{\fM}{\mathfrak{M}}
\newcommand{\fI}{\mathfrak{I}}
\newcommand{\fU}{\mathfrak{U}}
\newcommand{\fv}{\mathfrak{v}}
\newcommand{\fX}{\mathfrak{X}}
\newcommand{\fW}{\mathfrak{W}}
\newcommand{\fQ}{\mathfrak{Q}}
\newcommand{\fP}{\mathfrak{P}}
\newcommand{\fZ}{\mathfrak{Z}}
\newcommand{\one}{\ensuremath{\mathord \ast}}
\newcommand{\two}{\ensuremath{\mathord \dagger}}
\newcommand{\CMdef}{\operatorname{CMdef}}
\newcommand{\fC}{\mathfrak{C}}
\newcommand{\uc}{\underline{c}}

% Diagram helper
\def\mapdown#1{\Big\downarrow\rlap{$\vcenter{\hbox{$\scriptstyle#1$}}$}}

% Line spacing
\renewcommand{\baselinestretch}{1.27}

% Proof environment
\newenvironment{prf}[1][Proof]{\begin{proof}[\bf #1]}{\end{proof}}

\begin{document}
	
	\author[M. Asgharzadeh ]{Mohsen Asgharzadeh }
	\date{}
	\title[Perfect closure detects injective dimension ]{Perfect closure detects injective dimension  }
	
	\address{M. Asgharzadeh}
	\email{mohsenasgharzadeh@gmail.com}
	
 
	\subjclass[2020]{13D07; 13A35; 13B22}
	\keywords{absolute integral closure perfect closure; injective dimension; projective dimension; Gorenstein injective dimension.}
	
	\begin{abstract}
Let $R$ be a noetherian local ring of prime characteristic $p$, and let $R^\infty$ denote its perfect closure. We prove that a finitely generated \(R\)-module $N$ has finite injective dimension if and only if $\operatorname{Ext}_R^i(R^\infty, N) = 0$ for all $i > 0$. As a Gorenstein counterpart, we show that finiteness of the Gorenstein injective (or projective) dimension of $R^\infty$ forces $R$ to be Gorenstein, and we relate this to a non-noetherian Cohen factorization of $R \to R^\infty$. Applications include preservation of Gorensteinness under weakly  etale extensions, structural results on F-coherent and weakly F-nilpotent rings, along with the ascent of Frobenius closure of a parameter ideal  to its powers. Finally,  assuming
$\operatorname{Ext}^{i}_{R}(\frac{R}{\mathfrak m},R^{\infty})=0$
for some $i>\dim(R)$, we show $R$ is regular. This  has some applications.
	\end{abstract}

\maketitle
\tableofcontents
	

\section{Introduction}

A classical theme in commutative algebra is to detect homological invariants of modules over Noetherian rings through vanishing conditions. In characteristic \(p > 0\), the Frobenius map and its iterates provide a powerful tool for such investigations. Let \(R\) be a Noetherian local ring of characteristic \(p\), and let \(\mathrm{F} : R \to R\) denote the Frobenius map given by \(a \mapsto a^p\). For each integer \(e \ge 0\), we write \({}^e R\) for the \(R\)-algebra whose structure is induced by the \(e\)-th iterate \(\mathrm{F}^e : R \to R\).

In a foundational work, Herzog \cite{her} established the following celebrated criterion:

\begin{theorem}[Herzog \cite{her}]
	\label{thm:herzog}
	Let \(R\) be an \(F\)-finite local ring of characteristic \(p\), and let \(M\) be a finitely generated \(R\)-module. If \(\operatorname{Ext}_R^i({}^e R, M) = 0\) for all \(i > 0\) and infinitely many \(e\), then \(\operatorname{id}_R M < \infty\).
\end{theorem}

While elegant, Herzog's theorem requires checking infinitely many modules \({}^e R\)—one for each Frobenius iterate—to conclude finite injective dimension. This raises a natural question: can the entire family \(\{{}^e R\}_{e \ge 0}\) be encoded into a single test module that simultaneously detects finite injective dimension?
The perfect closure of \(R\), defined as the direct limit

\[
R^\infty := \varinjlim \bigl( R \xrightarrow{\mathrm{F}} R \xrightarrow{\mathrm{F}} R \xrightarrow{\mathrm{F}} \cdots \bigr)
= \varinjlim_{e \ge 0} ({}^e R),
\]

is a natural candidate. This \(R\)-algebra is the colimit of all Frobenius twists and has remarkable homological properties. Bhatt and Scholze \cite{BS1} showed that \(R^\infty\) has finite global dimension in a suitable sense, while   \cite{moh} established that every \(R^+\)-module has finite flat dimension over \(R^+\). Here, \(R^+\) denotes the absolute integral closure of an integral domain \(R\), defined as the integral closure of \(R\) inside an algebraic closure of its fraction field.

The main purpose of this note is to prove that the perfect closure \(R^\infty\) indeed serves as a universal test module for finite injective and projective dimensions. Our first main result is the following:

\begin{theorem}
	\label{thm:main1}
	Let \(R\) be a complete local ring of prime characteristic, and let \(N\) be a finitely generated \(R\)-module. Then the following are equivalent:
	\begin{enumerate}
		\item \(\operatorname{Ext}_R^i(R^\infty, N) = 0\) for all \(i > 0\),
		\item \(\operatorname{id}_R(N) < \infty\).
	\end{enumerate}
\end{theorem}
It maybe worth passing to the direct limit is not formally automatic as $\Ext(-,N)$ is not of this style. Indeed, the proof uses some spectral sequence
arguments together with the result of Bhatt-Scholze.
This result has the following corollaries:

\begin{corollary}
	\label{cor:tor}
	Let \((R,\mathfrak m,k)\) be a local ring of prime characteristic, and let \(L\) be an \(\mathfrak m\)-torsion \(R\)-module. Suppose that
	\(
	\operatorname{Tor}_i^R(R^\infty, L) = 0
	\)
	for all \(i > 0\). Then \(\operatorname{pd}_R(L) < \infty\), and consequently \(R\) is Cohen--Macaulay.
\end{corollary}

\begin{corollary}
	\label{cor:gorenstein1}
	Let \(R\) be a local ring of prime characteristic. Then \(R\) is Gorenstein if and only if \(\operatorname{Ext}_R^i(R^\infty, R) = 0\) for all \(i > 0\).
\end{corollary}

\begin{corollary}
	\label{cor:gorenstein2}
	Let \(R\) be a complete local ring of prime characteristic. Then \(R\) is Gorenstein if \(R^\infty\) is of finite Gorenstein projective dimension.
\end{corollary}

In order to connect the Gorenstein analogue of Theorem \ref{thm:main1}, our next results center on  interconnected themes:
\begin{enumerate}
	\item detecting Gorensteinness via the Gorenstein injective dimension of \(R^\infty\);
	\item structural consequences of \(F\)-coherent and weak \(F\)-nilpotent;

\end{enumerate}

\medskip

\noindent
(1): Detection of Gorensteinness via two approaches. 
We prove the following detection theorem:

\begin{theorem}
	\label{thm:gid}
	Let \((R,\mathfrak m)\) be a complete local ring of prime characteristic. If
	\(
	\operatorname{Gid}_R(R^\infty) < \infty,
	\)
	then \(R\) is Gorenstein.
\end{theorem}

As an application, we show the following:

\begin{corollary}
	\label{cor:weaklyetale}
	Suppose \(R \to S\) is weakly \'etale, \(S\) is complete, and \(R\) is Gorenstein. Then \(S\) is Gorenstein.
\end{corollary}



Section 3 suggests second, independent approach to the same detection problem—namely, detecting Gorensteinness from \(\operatorname{Gid}_R(R^\infty)\)—is provided by Cohen factorization. For a complete local domain \(R\), we construct a flat extension \(R \to Q = \varinjlim_e Q_e\), where each \(Q_e\) is noetherian and the transition maps are flat when \(R\) contains a regular purely inseparable subextension. This yields a Cohen factorization
\(
R \xrightarrow{\text{flat}} Q \twoheadrightarrow R^\infty,
\) with \(Q\) coherent. The key open question (see Question~\ref{25}) asks whether \(\operatorname{Gid}_Q(Q/I) < \infty\) implies \(\operatorname{id}_Q(Q) < \infty\). A positive answer would imply that \(R\) is Gorenstein via flat descent. Thus, the Cohen factorization approach offers a potentially powerful alternative route that is of independent interest.

\medskip


(2): The main focus of Section~5 is the following question from \cite{Al}:

\begin{quote}
	\emph{If every parameter ideal of \(R\) is Frobenius closed, does it follow that every power of each parameter ideal is Frobenius closed?}
\end{quote}

We partially answer this affirmatively in Proposition \ref{5.4}. This exemplifies a recurring theme of the paper: the perfection \(R^\infty\) is not merely an auxiliary construction but a powerful test object that encodes and transmits homological information.
We also show in Proposition \ref{4.5} that every complete \(F\)-coherent reduced ring is a quotient of a regular ring by a cohomologically complete-intersection ideal. 
We further establish several properties of weak \(F\)-nilpotent and \(F\)-coherent.



In Section 6 we resolve a conjecture \cite[Conjecture 5.19]{Elham} by showing that:
\begin{theorem}\label{18}
	Let $(R,\mathfrak m)$ be a local ring of characteristic $p>0$. 
	Assume that
	\(
	\operatorname{Ext}^{i}_{R}(\frac{R}{\mathfrak m},R^{\infty})=0
	\)
	for some $i>\dim(R)$. Then $R$ is regular.
\end{theorem}


Theorem \ref{18} has the following two applications:

\begin{enumerate}
	\item[(i)]
	\(
	\operatorname{Ext}^{+}_{R}(R^\infty,R^\infty)=0 \Longrightarrow 
	\text{$R$ is regular}.
	\)
	
	\item[(ii)]
	\(
	[\operatorname{Tor}^{R}_{i}(k,R^\infty)=0 \text{ for some } i>0] \Longrightarrow 
	\text{$R$ is regular}.
	\)
\end{enumerate}

The first drops the isolated singularity hypothesis from \cite{Ryo} when working with the minimal perfect closure. 
The second recovers a result of Aberbach--Li \cite{Ab} via a short and novel argument.

 Finally, we suggest that interested readers consult the forthcoming paper \cite{Ryo} for further homological properties of rings and modules in the context of perfect algebras, as well as for mixed-characteristic analogues of the results presented here.
	\maketitle
		\section{Perfect closure detects finite homological dimensions}
	
	
	
Let $(R,\mathfrak m,k)$ be a Noetherian local ring of  prime characteristic $p$.
The notation \(\pd_R(-)\) (resp. \(\id_R(-)\)) stands for the projective (resp. injective) dimension of \((-)\).

\begin{proposition}
	Let $R$ be a local ring of prime characteristic, and let
	$N$ be  either finitely generated  or a complete $R$-module. Then the following are equivalent:
	
	\begin{enumerate}
		\item
$
		\operatorname{Ext}_R^i(R^\infty,N)=0
		\qquad\text{for all } i>0,
	$
		
		\item
$
		\operatorname{id}_R(N)<\infty.
$
	\end{enumerate}
\end{proposition}

\begin{proof} We present the proof when $M$ is finitely generated.
	
	\noindent
	$(2)\Rightarrow(1)$.
	It follows by the so called Bass' conjecture (see \cite[9.6.2]{BH}), and   from $
	\operatorname{id}_R(N)<\infty
	$ that $R$ is Cohen--Macaulay. Let $\underline{x}:=x_1,\ldots, x_d$ be a parameter sequence. Then $\underline{x}^{[q]}:=x_1^{[q]},\ldots, x_d^{[q]}$ is an $R$-regular sequence. In other words, $\underline{x}$ is $R^{1/q}$-regular, i.e., ${}^e R$ is Cohen-Macaulay. Taking direct limit,
  $R^\infty$ is a balanced big Cohen--Macaulay module.
	By \cite[Corollary~7.7]{sim} we obtain the claim.
Here is a more elementary argument.
	Any module of finite injective dimension,  has finite $\omega$-resolution (see \cite[3.3.28(b)]{BH}), hence we may assume
	that
$
	N=\omega_R .
$
Recall that $\underline{x}$ is a parameter sequence.	Then the exact sequence
	\[
	0\longrightarrow
	R 
	\longrightarrow \oplus R_{x_i}
\longrightarrow
	\cdots
	\longrightarrow
	R_{x_1\ldots x_d}
	\longrightarrow
	\omega_R^\vee
	\longrightarrow
	0
	\]
	remains exact after tensoring with $R^\infty$, since $R^\infty$ is balanced
	big Cohen--Macaulay and the homologies are just the local cohomology of $R^\infty$. Hence
$
	\operatorname{Tor}_i^R(R^\infty,\omega_R^\vee)=0.
$
	Therefore, taking a Matlis dual, and using second isomorphism in  Fact \ref{fact:matlis-duality}, we have
$
	\operatorname{Ext}_R^i(R^\infty,\omega_R)	=\operatorname{Ext}_R^i(R^\infty,\omega_R^{\vee\vee})=0.
$
	

	
	\noindent
	$(1)\Rightarrow(2)$.
	There is a spectral sequence (see \cite[Theorem 10.74]{Rot}):
	\[
	E_2^{pq}
	=
	\operatorname{Ext}_{R^\infty}^{\,p}
	\!\left(
	L,
	\operatorname{Ext}_R^{\,q}(R^\infty,N)
	\right)
	\Longrightarrow
	\operatorname{Ext}_R^{\,p+q}(L,N),
	\]
	where $L$ is any $R^\infty$-module.
	By assumption (1), the spectral sequence collapses at $q=0$.
	Hence, we lead to the following isomorphism
	\[
	\operatorname{Ext}_{R^\infty}^{\,p}
	(L,\operatorname{Hom}_R(R^\infty,N))
	\cong
	\operatorname{Ext}_R^{\,p}(L,N)
\quad(\ast)	\]
	for every $p$.
		Recall from  Bhatt--Scholze \cite{BS1} that
	\[
	g:=\operatorname{gldim}(R^\infty)<\infty.
	\]

Set $L:=R^\infty/\fm R^\infty$. Then as an $R$-module, we have $L=\oplus R /\fm $, since $\fm L=0$.
	Therefore, the left hand side of $ (\ast)$ is zero for all $p>g$. It follows
	\[
	\operatorname{Ext}_R^{\,p}(L,N)=	\operatorname{Ext}_R^{\,p}(\oplus R /\fm ,N)=\prod \operatorname{Ext}_R^{\,p}(  R /\fm ,N)=0
	\qquad\text{for } p\gg0.
	\] Thus, 	$
  \operatorname{Ext}_R^{\,p}(  R /\fm ,N)=0
$  {for all} $p\gg0$.
 	By Roberts' theorem (see \cite[Theorem 2]{Ro2}), $\mu^n(\mathfrak m,N)$ is nonzero only in the range
$
	\bigl[\operatorname{depth}(N),\,\operatorname{id}(N)\bigr].
$
	Hence
$
	\operatorname{id}(N)<\infty.
$
\end{proof}
	
Let $(-)^\vee=\operatorname{Hom}_R(-,E_R(k))$ denote the Matlis duality
functor.

\begin{fact}(See \cite[4.1]{sim}).
	\label{fact:matlis-duality}
	Let $L$, $L'$, and $N$ be $R$-modules such that $N$ is Noetherian. Since $E_R(k)$ is injective we have following isomorphisms:
	\[
	\operatorname{Ext}^i_R(N, L')^\vee \cong \operatorname{Tor}^R_i(N, L'^\vee), \qquad
	\operatorname{Tor}^R_i(L, L')^\vee \cong \operatorname{Ext}^i_R(L, L'^\vee),
	\]where the second isomorphism is a consequence of Hom–tensor adjointness.
\end{fact}	
	
	
	
\begin{corollary}
	Let $(R,\mathfrak m,k)$ be a local ring of prime characteristic, and let
	$L$ be an $\fm$-torsion $R$-module. Suppose that
$
	\operatorname{Tor}_i^R(R^\infty,L)=0$ 
	 {for all positive } $i$.
	Then
$
	\operatorname{pd}_R(L)<\infty,
$
	and consequently $R$ is Cohen--Macaulay.
\end{corollary}

\begin{proof}
In view of  second isomorphism in  Fact \ref{fact:matlis-duality}, we obtain
	\[
	0=
	\operatorname{Tor}_i^R(R^\infty,L)^\vee
	\cong
	\operatorname{Ext}_R^i(R^\infty,L^\vee).
	\]
	
	Since $L$ is $\fm$-torsion, and by \cite[\S 4.2, Lemma]{sim} $L^\vee$ is complete. Hence, by the previous
	proposition,
$
	\operatorname{id}_R(L^\vee)<\infty.
$
Again by previous fact,
$
	\operatorname{fd}_R(L)<\infty,
$
	and therefore
$
	\operatorname{pd}_R(L)<\infty.
$	Finally, the existence of a nonzero $\fm$-torsion module of finite projective
	dimension implies that $R$ is Cohen--Macaulay (see e.g. \cite[Observation 3.15]{moh5}).
\end{proof}
\begin{corollary}\label{25}
	Let $R$ be a local ring of prime characteristic. Then $R$ is Gorenstein if and only if $\operatorname{Ext}_R^+(R^\infty, R) = 0$.
\end{corollary}
In what follows, we assume some familiarity with Gorenstein homological algebra, see \cite{Totally}
for a nice survey.
The notation \(\Gpd_R(-)\) (resp. \(\Gid_R(-)\)) stands for the Gorenstein  projective (resp. injective) dimension of \((-)\).
\begin{corollary}
	Let $R$ be a complete local ring of prime characteristic. If $R^\infty$ has finite Gorenstein projective dimension, then $R$ is Gorenstein.
\end{corollary}

\begin{proof}
	Recall that $\operatorname{Gfd}_R(-) \leq \operatorname{Gpd}_R(-)$, and that for every $R$-module $M$ of finite Gorenstein flat dimension, we have the equality
	\[
	\operatorname{Gfd}_R(M) = \sup \left\{ \operatorname{edepth}(R_{\mathfrak p}) - \operatorname{edepth}_{R_{\mathfrak p}}(M_{\mathfrak p}) \mid \mathfrak p \in \operatorname{Spec} R \right\},
	\]where $\operatorname{edepth}$ means $\Ext-depth$.
	Applying this to $M = R^\infty$, and using the fact that $(R^\infty)_{\mathfrak p} \cong (R_{\mathfrak p})^\infty$, we obtain
	\[
	\operatorname{edepth}_{R_{\mathfrak p}}((R^\infty)_{\mathfrak p}) 
	\geq \operatorname{depth}_{R_{\mathfrak p}}((R^\infty)_{\mathfrak p}) 
	\geq \operatorname{depth}_{R_{\mathfrak p}}(R_{\mathfrak p}) 
	= \operatorname{edepth}_{R_{\mathfrak p}}(R_{\mathfrak p}).
	\]
	Hence $\operatorname{Gfd}_R(R^\infty) \leq 0$.
	Moreover, if $\operatorname{Gfd}_R(M)$ is finite, then
	\[
	\operatorname{Gfd}_R(M) = \sup \left\{ i \in \mathbb{Z} \mid \operatorname{Tor}_i^R(E, M) \neq 0 \text{ for some injective } R\text{-module } E \right\}.
	\]
	In particular, $\operatorname{Tor}_i^R(E_R(k), R^\infty) = 0$ for all $i > 0$. Taking the Matlis dual and applying the second isomorphism in Fact \ref{fact:matlis-duality}, we get
	\[
	\operatorname{Ext}_R^i(R^\infty, E_R(k)^\vee) 
	= \operatorname{Ext}_R^i(R^\infty, R) = 0.
	\]
	Then, by Corollary \ref{25}, we conclude that $R$ is Gorenstein.
\end{proof}


	When is $\operatorname{Gid}_R(R^\infty) < \infty$?
A natural answer is precisely in the Gorenstein case. Note that $\operatorname{Gid}_R(R^\infty) \leq \operatorname{id}_R(R^\infty)$, with equality whenever $R^\infty$ has finite injective dimension over $R$ (see \cite[Proposition 3.10]{Totally}). The latter property has been fully investigated   in \cite{Ryo}.
 
 \begin{remark}
Suppose $R$ is F-splitting. Then $\operatorname{Gid}_R(R^\infty) < \infty$ if and only if $R$ is Gorenstein. 
 \end{remark}


\begin{proof}Suppose $\operatorname{Gid}_R(R^\infty) < \infty$.
Let $E$ be an injective $R$-module. There is an embedding $\Ext^j_R(E,R)\subseteq   \Ext^j_R(E,R^\infty)=0$ for all $j\gg 0$ (see \cite[Corollary 3.7]{Totally}). Then, $\Ext^j_R(E,R)=0$ for all $j\gg 0$, and so $\operatorname{Gid}_R(R) < \infty$. Since $R$ is projective this implies that $\operatorname{Gid}_R(R)=\id_R(R)$ (see \cite[Proposition 3.11]{Totally}), hence
 $R$ is Gorenstein.
\end{proof}






Against to $\Tor(E,-)$, the functor $\Ext(E,-)$ is not behaved with   respect to the flat base change. This is the main difficulty in working with $\operatorname{Gid}_R(R^\infty)$.
In the next sections we present a systematic study of Gorenstein injective dimension of $R^\infty$.

	\maketitle
	
\section{Cohen's factorization for perfect closure}

Let $k$ be a perfect field, and let $(R,\mathfrak m,k)$ be a Noetherian complete  local domain of
characteristic $p>0$. For each $e\geq 0$, consider the local homomorphism
$
\varphi_e:R\to R^{1/p^e}.
$
Choose a Cohen factorization
\(
R\xrightarrow{\dot\varphi_e}Q_e
\stackrel{\varphi'_e}\twoheadrightarrow\widehat{R^{1/p^e}}={R^{1/p^e}}.\)
Thus $Q_e$ is a complete Noetherian local ring,
$\dot\varphi_e$ is flat, and the closed fiber
$
Q_e/\mathfrak mQ_e
$
is regular.
By the weak functoriality of Cohen's factorizations, the
factorizations may, after suitable choices, be made compatible with
the maps
$
R^{1/p^e}\to R^{1/p^{e+1}}.
$
In particular, we obtain a directed system
\[
Q_0\longrightarrow Q_1\longrightarrow Q_2
\longrightarrow\cdots.
\]
We emphasize that the transition maps in above are not asserted to
be flat, but see {Lemma} \ref {cf} below.  The existence of compatible maps
between suitable Cohen factorizations is the weak functoriality
property established by Nasseh and Sather-Wagstaff \cite{sa}.


\begin{lemma}
	With notation as above, the canonical homomorphism
	$
	R\to  Q=\varinjlim_eQ_e
	$
	is flat.
\end{lemma}

\begin{proof}
	Let
	$
	\zeta:=	0\to M'\to M\to M''\to0
	$
	be an exact sequence of $R$-modules. Since $R\to Q_e$ is flat for
	every $e$, tensoring with $Q_e$ gives an exact sequence
	$
	\zeta\otimes_RQ_e
	.
	$
	Hence
	$
	\varinjlim_e(\zeta\otimes_RQ_e)	
	$
	is exact. Since tensor product commutes with filtered colimits,
	$
	\varinjlim_e(M\otimes_RQ_e)
	\cong
	M\otimes_R\left(\varinjlim_eQ_e\right)
	=M\otimes_RQ,
	$
	and similarly for $M'$ and $M''$. Consequently,
	\(
	\zeta\otimes_RQ
	\)
	is exact. 
\end{proof}

The ring $Q$ is generally non-Noetherian. Nevertheless, it is built
from Noetherian local rings, and this allows one to approximate
finitely presented $Q$-modules by modules over the $Q_e$.
We record \cite{G1} as a standard reference for coherent rings.
\begin{lemma}
	Assume that $Q$ is coherent. Let $N$ be a cyclic finitely presented
	$Q$-module. Then, for some $e$, there exists a cyclic finitely
	presented $Q_e$-module $N_e$ such that
	$
	N\cong N_e\otimes_{Q_e}Q.
	$
\end{lemma}

\begin{proof}
	Since $N$ is cyclic, write
	$
	N\cong Q/I
	$
	for an ideal $I$ of $Q$. Since $N$ is finitely presented and $Q$ is
	coherent, the ideal $I$ is finitely generated. Thus
	$
	I=(x_1,\ldots,x_r).
	$
	As
	$
	Q=\varinjlim_eQ_e,
	$
	each $x_i$ is represented by an element of $Q_e$ for all sufficiently
	large $e$. After increasing $e$, choose elements
	$
	x_{1,e},\ldots,x_{r,e}\in Q_e
	$
	whose images in $Q$ are $x_1,\ldots,x_r$. Put
	$
	I_e=(x_{1,e},\ldots,x_{r,e})\subseteq Q_e
	$
	and
	$
	N_e=Q_e/I_e.
	$
	Then
	$
	I_eQ=I,
	$
	and hence
	$
	N_e\otimes_{Q_e}Q
	\cong Q/I_eQ
	\cong Q/I
	=N.
	$
	Since $Q_e$ is Noetherian, $N_e$ is finitely presented. It is
	obviously cyclic.
\end{proof}

Set
\(
A=\bigcup_{e\geq 0}R^{1/p^e}.
\) We summarize the strategy in the following diagram:
\[
\begin{CD}
R @>>> Q=\varinjlim_eQ_e @> 
\twoheadrightarrow	>> A\\
@. @AAA @AAA\\
@. Q_e @>\twoheadrightarrow>> R^{1/p^e}
\end{CD}
\]
with
$
R\to Q_e$ {flat with regular closed fiber},
and
$
R\to Q
$  flat.
Assume that $B\subseteq R$ is a purely inseparable extension of
Noetherian domains of characteristic $p>0$, and that $B$ is regular.
The first observation is that $R$ and $B$ have the same perfect
closure.

\begin{fact}
	If $B\subseteq R$ is purely inseparable, then, inside a common
	algebraic closure of their fraction fields,
	$
	R^\infty=B^\infty.
	$ Indeed, since $B\subseteq R$, we have
	$
	B^\infty\subseteq R^\infty.
	$
	Conversely, let $x\in R$. Since $R/B$ is purely inseparable, there
	exists $e\geq 0$ such that
	$
	x^{p^e}\in B.
	$
	Hence
	$
	x\in B^{1/p^e}\subseteq B^\infty.
	$
	Thus $R\subseteq B^\infty$, and therefore
	$
	R^\infty\subseteq B^\infty.
	$
	The two inclusions give the assertion.
\end{fact}

 

Thus, in the purely inseparable situation, we may write
$
A=R^\infty=B^\infty.
$
For each $e\geq0$, consider
$
C_e:=R\otimes_B B^{1/p^e}.
$
There are canonical maps
$
C_e\to C_{e+1}
$
induced by
$
B^{1/p^e}\to B^{1/p^{e+1}}.
$ The regularity of $B$ now becomes crucial.

\begin{lemma}\label{cf}
	If $B$ is regular, then the homomorphism
	\(
	B^{1/p^e}\to B^{1/p^{e+1}}
	\)
	is flat for every $e\geq0$. Consequently,
	\(
	C_e\to C_{e+1}
	\)
	is flat. In particular, we get a Cohen factorizations with coherent rings $\{Q,A\}$:
	\[
	R\stackrel{flat}\longrightarrow Q\twoheadrightarrow A:={R^{\infty}},
	\]
\end{lemma}

\begin{proof}
	The map
	\(
	B^{1/p^e}\to B^{1/p^{e+1}}
	\)
	is obtained from the Frobenius homomorphism on $B^{1/p^e}$.
	Since $B$ is regular, every iterate of Frobenius is flat by Kunz's
	theorem. Hence
	\(
	B^{1/p^e}\to B^{1/p^{e+1}}
	\)
	is flat.
	Moreover,
	\[
	C_{e+1}
	=
	R\otimes_BB^{1/p^{e+1}}
	\cong
	\left(R\otimes_BB^{1/p^e}\right)
	\otimes_{B^{1/p^e}}B^{1/p^{e+1}}
	=
	C_e\otimes_{B^{1/p^e}}B^{1/p^{e+1}}.
	\]
	Thus $C_e\to C_{e+1}$ is a base change of the flat map
	\(
	B^{1/p^e}\to B^{1/p^{e+1}},
	\)
	and hence is flat.
\end{proof}


\begin{question}\label{25}
	If
	\(
	\operatorname{Gid}_Q(Q/I)<\infty ,
	\)
	then is $\id_Q(Q)<\infty$?	\end{question}
The point is that $Q$ is not noetherian (see \cite[Theorem 3.22]{Totally}).	Since $R\to Q$ is flat, positive answer to {Question} \ref{25} implies $R$  is Gorenstein.
\maketitle


\section{Gorenstein injective dimension detected by perfect closure}


Shimomoto \cite{Shi10} defines \(R\) to be \emph{F-coherent} if the perfect closure \(R^\infty\) is coherent.

\begin{theorem}\label{3.1}
	Suppose \((R,\mathfrak m)\) is an \(F\)-coherent ring. If
	\(
	\operatorname{Gid}_R(R^\infty) < \infty,
	\)
	then \(R\) is Gorenstein.
\end{theorem}

\begin{proof}
	Recall that any finitely generated module over a coherent ring has a resolution by finitely generated free modules, and in fact over \(R^\infty\) such a resolution is of finite length (see \cite{A}). In particular, we have an \(R^\infty\)-resolution:
	\[
	0 \longrightarrow (R^\infty)^{n_t} \longrightarrow \cdots \longrightarrow (R^\infty)^{n_0} \longrightarrow R^\infty / \mathfrak m R^\infty \longrightarrow 0,
	\]
	where the \(n_i\) are integers. Since \(R \subseteq R^\infty\), the above exact sequence is indeed \(R\)-linear.
	
	We use the following two properties of Gorenstein injective modules (see \cite[Proposition 3.9, Proposition 3.3]{Totally}):
	
	\begin{enumerate}
		\item Finite direct sums of Gorenstein injective modules are Gorenstein injective.
		\item If
		\(
		0 \longrightarrow M_1 \longrightarrow M_2 \longrightarrow M_3 \longrightarrow 0
		\)
		is a short exact sequence and any two of the modules have finite Gorenstein dimension, then so does the third.
	\end{enumerate}
	
	Property (1) implies that any finite direct sum of modules of finite injective dimension has finite Gorenstein injective dimension.
	
	Since \(R^\infty\) has finite Gorenstein injective dimension, the above resolution implies that
	\(
	R^\infty / \mathfrak m R^\infty
	\)
	has finite Gorenstein dimension. Now \(R / \mathfrak m\) is a direct summand of
	\(
	R^\infty / \mathfrak m R^\infty.
	\)
	It follows from \cite[Proposition 3.3]{Totally} that
	\[
	\operatorname{Gid}_R(R / \mathfrak m) < \infty.
	\]
	Therefore, \(R\) is Gorenstein (see \cite[Corollary 3.23]{Totally}).
\end{proof}

The following is the main result of this section.

\begin{theorem}\label{4.2}
	Suppose \((R,\mathfrak m)\) is a complete local ring. Then
	\(
	\operatorname{Gid}_R(R^\infty) < \infty
	\)
	if and only if  \(R\) is Gorenstein.
\end{theorem}

\begin{proof}
	Since \(R\) is a homomorphic image of a Gorenstein ring of finite Krull dimension, the class of Gorenstein injective modules is closed under direct limits (see \cite[Theorem 3.16]{Totally}). This enriches property (1) from Theorem \ref{3.1}. Recall from \cite{BH} that the global dimension of \(R^\infty\) is finite, and use the argument of Theorem \ref{3.1} to deduce that \(R\) is Gorenstein. The reverse implication is true without any assumption, and indeed it is well-known (see \cite[Corollary 3.23]{Totally}).
\end{proof}
The following gives a new proof of  \cite[Example 3.17(1)]{Ryo}.

\begin{corollary}\label{4.3}
	Suppose 
	\(
	\operatorname{id}_R(R^\infty) < \infty 	\) (resp. 	\(
	\operatorname{fdim}_R(R^\infty) < \infty
	\)).
Then \(R\) is regular.
\end{corollary}

\begin{proof}
By the argument of Theorem \ref{3.1}, it turns out that   
\(
\operatorname{id}_R(k) < \infty	\) (resp. 	\(
\operatorname{fdim}_R(k) < \infty
\)). 
So, \(R\) is regular.
\end{proof}

Suppose
$
\operatorname{Gid}_R(M \otimes R^\infty) < \infty.
$
One may ask whether \(\operatorname{id}_R(M) < \infty\). This is not the case, as the next result indicates. 

\begin{example}
Let $(R,\fm)$ be a  Gorenstein ring which is not  integral domain, and take \(M := \mathfrak p\) where \(\mathfrak p \in \operatorname{Spec}(R)\). Clearly,
	\(
	\operatorname{Gid}_R(M \otimes R^\infty) < \infty,
	\) as $R$ is  Gorenstein.
	If \(\operatorname{id}(\mathfrak p) < \infty\), then we would have \(\operatorname{pd}(\mathfrak p) < \infty\), as \(R\) is Gorenstein. This would imply that \(R\) is a domain, a contradiction.
\end{example}

An extension \(A \to B\) is called \emph{weakly \'etale} (or \emph{absolutely flat}) if both \(A \to B\) and \(B \otimes_A B \to B\) are flat.

\begin{corollary}\label{4.4}
	Suppose \(R \to S\) is weakly \'etale, \(S\) is complete, and \(R\) is Gorenstein. Then \(S\) is Gorenstein.
\end{corollary}

\begin{proof}
	Since \(R \to S\) is weakly \'etale, we have \(S^{1/p^n} \simeq S \otimes_R R^{1/p^n}\). In particular,
	\(
	S^\infty \simeq S \otimes_R R^\infty.
	\)
	Since \(R\) is Gorenstein, Theorem \ref{4.2} gives \(\operatorname{Gid}_R(R^\infty) < \infty\). This implies
	\[
	\operatorname{Gid}_S(S \otimes_R R^\infty) < \infty,
	\]
	and hence \(\operatorname{Gid}_S(S^\infty) < \infty\). By Theorem \ref{4.2}, \(S\) is Gorenstein.
\end{proof}

\begin{question}
	\begin{enumerate}
		\item[1)] (See \cite{Ryo}).	Suppose \(
		\operatorname{Ext}^{+}_{R}(R^\infty,R^\infty)=0\).
		Is 	$R$ regular?
		
		\item[2)](See \cite[Conjecture 5.19]{Elham}). Suppose \(
		\operatorname{Ext}^{\gg 0}_{R}(k,R^\infty)=0\).
		Is 	$R$ regular? 
		
	\end{enumerate}
\end{question}



In Section 6 we need the following:

\begin{proposition}\label{47} Adopt the above notation. Any positive answer to
	Question	2),  yields a positive answer to Question 1). In particular, 1) is valid over F-pure rings.
\end{proposition}
\begin{proof}
	There is a spectral sequence
	\[
	\operatorname{Ext}^{i}_{R^\infty}
	\left(
	L,\operatorname{Ext}^{j}_{R}(R^\infty,R^\infty)
	\right)
	\Longrightarrow
	\operatorname{Ext}^{i+j}_{R}(L,R^\infty),
	\]where $L$ is any $R^\infty$-module.
	This gives 
	$\operatorname{Ext}^{i}_{R}
	\left(
	L,\operatorname{Hom}_{R}(R^\infty,R^\infty)
	\right)
	\simeq
	\operatorname{Ext}^{i}_{R}(L,R^\infty).
	$ Set
	$
	L=\frac{R^\infty}{\mathfrak m R^\infty}.
	$
	But
	$
	\operatorname{gl.dim}(R^\infty)<\infty,
	$ (see \cite{BS1}).	Thus
	$
	\operatorname{Ext}^{i}_{R}(L,-)=0 $
	for all $i\gg0,
	$
	and therefore
	$
	\operatorname{Ext}^{i}_{R}
	\left(
	\frac{R }{\mathfrak m },
	R^\infty
	\right)=0 
	$,  and so $R$ is regular.
\end{proof}	
	\maketitle
	
\section{Ascent and descent from (to) $R^\infty$ }

The following was asked in \cite[Question 3.9]{Al}:
\begin{question} Assume each parameter ideal in $R$ is Frobenius closed. Is it true that every power of a parameter ideal is also Frobenius closed?
\end{question}

\begin{notation}Let $I\lhd R$ be an ideal.
	By $I^\star$ (resp. $I^F$) we denote the tight closure (resp. Frobenius closure) of $I$.
\end{notation}
\begin{proposition}\label{42}
	Let $R$ be a equi-dimensional and reduced local ring of prime characteristic $p$ which is quotient of a Cohen-Macaulay ring.
	Assume that $R$ is $F$-coherent. Let
	$
	I=(g_1,\ldots,g_d)
	$
	be an ideal generated by a system of parameters of $R$. If $I$ is
	Frobenius closed, then
	$
	(I^n)^F=I^n
	$
	for every integer $n\geq 1$.
\end{proposition}
\begin{proof}
	By F-coherent we see from \cite[Proposition 3.11]{Shi10} that $I^\star=I^F=I$, and by equi-dimensional assumption it follows that $R$ is F-rational (see \cite[Proposition 10.3.5]{BH}). Then $R$ is Cohen-Macaulay  (see \cite[Corollary 10.3.4]{BH}), as it is quotient of a Cohen-Macaulay ring. So, $I$ is generated 	by regular sequence and the claim is clear by 
	\cite[Proposition 3.10]{Al}.
\end{proof}

Here, we are going to drop the restriction on the length of system of parameters, the equi-dimensional,  reduced, and quotient of a Cohen-Macaulay ring 
from {Proposition} \ref{42}. We record a well-known fact:
\begin{fact}[Frobenius descent]
	Let $I$ be an ideal of $R$. If
	$
	x\in IR^\infty\cap R,
	$
	then there exists an integer $e\geq 0$ such that
	$
	x^{p^e}\in I^{[p^e]}.
	$
In particular, $I^F=IR^\infty\cap R$.
\end{fact}
	
	\iffalse	\begin{proof}
		Since
		$
		x\in IR^\infty,
		$
		there exist elements
		$
		a_1,\ldots,a_s\in I
		$
		and
		$
		y_1,\ldots,y_s\in R^\infty
		$
		such that
		$
		x=\sum_{j=1}^s a_jy_j.
		$
		Since only finitely many elements $y_j$ occur and
		$
		R^\infty=\bigcup_{e\geq 0}R^{1/p^e},
		$
		there exists $e\geq 0$ such that
		$
		y_j\in R^{1/p^e}
		$
		for every $j$. Hence
		$
		y_j^{p^e}\in R
		$
		for every $j$.
		Raising the equality
		$
		x=\sum_{j=1}^s a_jy_j
		$
		to the $p^e$-th power gives
		$
		x^{p^e}
		\sum_{j=1}^s a_j^{p^e}y_j^{p^e},
		$
		because the Frobenius map is a ring homomorphism in characteristic $p$. Since
		$
		a_j\in I,
		$
		we have
		$
		a_j^{p^e}\in I^{[p^e]},
		$
		while
		$
		y_j^{p^e}\in R.
		$
		Therefore
		$
		x^{p^e}\in I^{[p^e]}.
		$
		Now suppose that $I$ is Frobenius closed. If
		$
		x\in IR^\infty\cap R,
		$
		then, by the first part,
		$
		x^{p^e}\in I^{[p^e]}
		$
		for some $e\geq 0$. Hence
		$
		x\in I^F=I.
		$
		Thus
		$
		IR^\infty\cap R\subseteq I.
		$
		The reverse inclusion is immediate, and consequently
		$
		IR^\infty\cap R=I.
		$
	\end{proof}\fi

Let us present a property weaker than F-coherent rings. Recall from \cite{MP} that $R$ is called \emph{weakly F-nilpotent}  if  every element of $H^i_\fm(R)$ is annihilated by some   $\F^n$ for
$i < \dim(R)$.
 Recall from \cite[Exercise 68]{MP}  that over excellent F-nilpotent rings also one has
$I^\star=I^F$, and so similar to F-coherent rings.
\begin{proposition}\label{5.4}
	Let $R$ be a  local ring of prime characteristic $p$.
	Assume that $R$ is weakly F-nilpotent. Let
	$
	I=(g_1,\ldots,g_t)
	$
	be an ideal generated by a part of a system of parameters of $R$. If $I$ is
	Frobenius closed, then
	$
	(I^n)^F=I^n
	$
	for every integer $n\geq 1$.
\end{proposition}
\begin{proof}
	We proceed by induction on $n$. The case $n=1$ follows from the hypothesis
	that $I$ is Frobenius closed.
	Since $R$ is weakly F-nilpotent,   $R^\infty$ is a
	balanced big Cohen--Macaulay $R$-algebra (see \cite[Proposition 12.22]{MP}). 
	Suppose that
	$
	(I^n)^F=I^n
	$
	for some $n\geq 1$, and let
	$
	r\in (I^{n+1})^F.
	$
	Since
	$
	I^{n+1}\subseteq I^n,
	$
	we have
	$
	r\in (I^n)^F=I^n.
	$
	Hence we can write
	$
	r=\sum_{|\alpha|=n}c_\alpha g^\alpha,
	$
	where
	$
	\alpha=(\alpha_1,\ldots,\alpha_t),
	$ and
	$|\alpha|=\alpha_1+\cdots+\alpha_t,
	$
	also
	$
	g^\alpha=g_1^{\alpha_1}\cdots g_t^{\alpha_t}.
	$
	Since
	$
	r\in (I^{n+1})^F,
	$
	there exists a power
	$
	q=p^e
	$
	of $p$ such that
	$
	r^q\in (I^{n+1})^{[q]}.
	$
	Set
	$
	h_i=g_i^q,1\leq i\leq t
	$
	and
	$
	Q=(h_1,\ldots,h_t)=I^{[q]}.
	$
	Since Frobenius powers commute with ordinary powers,
	$
	(I^{n+1})^{[q]}
	(I^{[q]})^{n+1}
	Q^{n+1}.
	$
	Therefore
	$
	r^q
	\sum_{|\alpha|=n}c_\alpha^q h^\alpha
	\in Q^{n+1},
	$
	where
	$
	h^\alpha=h_1^{\alpha_1}\cdots h_t^{\alpha_t}.
	$
	Because $R^\infty$ is balanced big Cohen--Macaulay and
	$g_1,\ldots,g_t$ is part of a system of parameters, the sequence
	$
	g_1,\ldots,g_t
	$
	is $R^\infty$-regular. Hence
	$
	h_1=g_1^q,\ldots,h_t=g_t^q
	$
	is also an $R^\infty$-regular sequence.
	Consider the homogeneous polynomial
	$
	F(Y_1,\ldots,Y_t)
	\sum_{|\alpha|=n}c_\alpha^qY^\alpha
	$
	in
	$
	R^\infty[Y_1,\ldots,Y_t].
	$
	The containment
	$
	F(h_1,\ldots,h_t)\in Q^{n+1}R^\infty
	$
	and the regularity of
	$
	h_1,\ldots,h_t
	$
	imply, by the coefficient criterion for powers of a regular sequence, that
	$
	c_\alpha^q\in QR^\infty
	$
	for every $\alpha$ with
	$
	|\alpha|=n.
	$
	Since
	$
	c_\alpha^q\in R,
	$
	the Frobenius descent lemma, applied to the ideal
	$
	Q=I^{[q]},
	$
	gives an integer $f\geq0$ such that
	$
	(c_\alpha^q)^{p^f}\in Q^{[p^f]}.
	$
	Thus
	$
	c_\alpha^{qp^f}
	\in
	(I^{[q]})^{[p^f]}
	I^{[qp^f]}.
	$
	Since
	$
	qp^f
	$
	is again a power of $p$, this shows that
	$
	c_\alpha\in I^F.
	$
	By the hypothesis
	$
	I^F=I,
	$
	we conclude that
	$
	c_\alpha\in I
	$
	for every $\alpha$. Consequently,
	$
	r
	=
	\sum_{|\alpha|=n}c_\alpha g^\alpha
	\in I^{n+1}.
	$
	Therefore
	$
	(I^{n+1})^F=I^{n+1}.
	$
	The induction is complete.
\end{proof}

There are several similarity between F-coherent and weakly nilpotent, namely the big Cohen-Macaulayness of their perfection.
\begin{proposition}
	Let $R$ be a reduced local ring of characteristic $p > 0$. Then $R$ is
	weakly nilpotent if and only if the henselization 
	is so.
\end{proposition}

\begin{proof}
This is parallel to \cite[Proposition 3.18]{Shi10}, so we leave its straightforward modification to the reader.
\end{proof}

An ideal $I$ of regular ring ring $R$ is called  cohomologically complete-intersection if
$$H^i_I(R)=0\Longleftrightarrow i\neq \grade(I,R).$$For example, any ideal generated by regular sequence is of this form. A celebrated theorem of Peskine-Szpiro \cite{PS} says any Cohen-Macaulay ideal in a regular ring of prime characteristic is of this mood. Here, is another sample:
\begin{proposition}\label{4.5} Any complete $F$-coherent reduced ring $A$ is the form
	\(
	A \cong   R/I,
	\)
	where $R$ is regular and $I$ is cohomologically complete-intersection.
\end{proposition}
\begin{proof} By Cohen's structure theorem, $A$ is the form
	\(
	A \cong   R/I,
	\)
	where $R$ is regular. Since $F$-coherent implies that $A^\infty$ is big Cohen--Macaulay (see \cite{Shi10}), we have
	\[
	H_\fm^i(A^\infty)=0 \qquad  \forall i\neq  \dim(A) ,
	\]Then we have
	\(
	H_I^i(R)=0\) for all  \(i\neq \operatorname{ht}(I),
	\) see  $(2)\equiv(6)$ in \cite[Example 4.38]{topo}.
	Therefore $A$ is a cohomologically complete-intersection ring.
\end{proof}


\begin{remark}\label{46}
	The converse of above result is not true. Indeed, there is a hypersurface $A$ which is not F-coherent (see \cite[Example 4.4]{Shi10}).
\end{remark}





We conclude this section with a characterization of weakly \(F\)-nilpotent complete local rings.

\begin{corollary}\label{59}
	A complete local ring is weakly \(F\)-nilpotent if and only if it is cohomologically complete-intersection.
\end{corollary}

\begin{proof}
	This follows from the proof of Proposition \ref{4.5}.
\end{proof}

We refer to Remark \ref{46} and Corollary \ref{59} to see an explicit difference between \(F\)-coherent rings and weakly \(F\)-nilpotent rings.
Here is another.

\begin{proposition}\label{obs:limits}
	Let \(\kappa\) be a limit ordinal such that \(\{(R_\gamma, \mathfrak m_\gamma)\}_{\gamma < \kappa}\) is a filtered direct system of weakly \(F\)-nilpotent rings, with \(\mathfrak m_\gamma R_\delta = \mathfrak m_\delta\) for all \(\gamma < \delta < \kappa\). Then \(\varinjlim_{\gamma} R_\gamma\) is weakly \(F\)-nilpotent.
\end{proposition}

\begin{proof}
	Let \(R := \varinjlim_{\gamma} R_\gamma\). By the main result of \cite{O}, \(R\) is Noetherian. Also, \(R\) is local with unique maximal ideal \(\mathfrak m = \varinjlim_{\gamma} \mathfrak m_\gamma\). Suppose \(\mathfrak m\) is generated by \(\underline{x} := x_1, \ldots, x_t\). Then there exists \(\lambda < \kappa\) such that \(x_i \in R_\gamma\) for all \(i\) and all \(\gamma \ge \lambda\). Moreover, \(\mathfrak m = \mathfrak m_\gamma R\) for all \(\gamma \ge \lambda\).
	In view of \cite[Theorem 15.1]{Mat}, we have
	\[
	\dim (R )\le \dim (R_\gamma) + \dim (R / \mathfrak m_\gamma R).
	\]
	Since \(\mathfrak m = \mathfrak m_\gamma R\), we have \(\dim (R / \mathfrak m_\gamma R) = 0\), and therefore \(\dim (R) \le \dim (R_\gamma)\).
	Let \(i < \dim R\). Now,
	\(
	R^\infty \cong \varinjlim_{\gamma < \kappa} (R_\gamma)^\infty,
	\)
	and so
	\[
	\begin{aligned}
	H_{\mathfrak m}^{i}(R^\infty)
	 \cong H_{\underline{x}R}^{i}(R^\infty) 
	&\cong H_{\underline{x}R}^{i}\left(\varinjlim_{\lambda \le \gamma < \kappa} (R_\gamma)^\infty\right) \\
	&\cong \varinjlim_{\lambda \le \gamma < \kappa} H_{\underline{x}}^{i}(R_\gamma^\infty)\cong \varinjlim_{\gamma < \kappa} H_{\mathfrak m_\gamma}^{i}(R_\gamma^\infty) = 0,
	\end{aligned}
	\]where $H_{\underline{x}}^{i}(-)$ stands for the $i$-th local cohomology with respect to $\underline{x}$.
	Thus \(R^\infty\) is big Cohen--Macaulay, and so \(R\) is weakly \(F\)-nilpotent (see \cite[Proposition 12.22]{MP}).
\end{proof}\begin{proposition}\label{obs:limits}
Let \(\kappa\) be a limit ordinal such that \(\{(R_i, \mathfrak m_i);f_{i,j}:R_i\to R_j\}_{i < \kappa}\) is a filtered direct system of  \(F\)-coherent rings, with \(\mathfrak m_\gamma R_\delta = \mathfrak m_\delta\) for all \(\gamma < \delta < \kappa\) and $f_{i,j}$ is  weakly \'etale. Then \(\varinjlim_{i} R_i\) is \(F\)-coherent.
\end{proposition}

\begin{proof}	Let \(R := \varinjlim_{i} R_\gamma\).
It follows by  change of coefficient (see \cite[page 46]{Mat}) along with the weakly \'etaleness that the following map \[f_{i,j}^\infty:(R_i)^\infty\stackrel{\cong}\lo R_i   \otimes_{R_i} {(R_i)^\infty}\stackrel{f_{i,j}\otimes1}\lo R_j\otimes_{R_i}(R_i)^\infty \stackrel{\cong}\lo(R_j)^\infty,\]
is flat. Since $(R)^\infty=\varinjlim_{i} \{(R_i)^\infty;f_{i,j}^\infty\}$
it follows from \cite[Theorem 2.3.3]{G1} that $R$ is \(F\)-coherent.
\end{proof}
\maketitle

\section{A conjecture on $\mu_i(R^\infty)=0$}


Recall that it was conjectured in  \cite[Conjecture 5.19]{Elham} that:

\begin{conjecture}\label{61}
Suppose \(
\operatorname{Ext}^{i }_{R}(k,R^\infty)=0\) for some $i>\dim(R)$.
Then	$R$ is  regular.

\end{conjecture} 
Mahdavi \cite{Elham} showed that the condition $i > \dim(R)$ is genuinely necessary, and established this over rings with isolated singularities. 

The following fact is well known.



\begin{fact}\label{comuf}(See \cite[Theorem~2.6.2]{G1}).
	Let $A$ be a commutative ring. Then following are equivalent:
	\begin{enumerate}
		\item[(i)]
		\(A \) is coherent;
		
		\item[(ii)] For all finitely presented  module $P$, any directed system $\{M_i\}$ of $A$-modules and all $n\geq 1$, we have $$\Ext^n_A(P,\vil M_i)\cong\vil\Ext^n_A(P,  M_i).$$
	\end{enumerate}
\end{fact}

In particular, Corollary \ref{4.3} extends as follows.

\begin{theorem}
	Let $(R,\mathfrak m)$ be a local ring of characteristic $p>0$. 
	Assume that
	\(
	\operatorname{Ext}^{i}_{R}\left(\frac{R}{\mathfrak m},R^{\infty}\right)=0
	\)
	for some $i>\dim(R)$. Then $R$ is regular.
\end{theorem}

\begin{proof}
Recall from \cite{BH} that the global dimension of \(R^\infty\) is finite.
In particular, we have an \(R^\infty\)-resolution:
\[
0 \longrightarrow \oplus_{X_t}  R^\infty   \stackrel{f_t}\longrightarrow \cdots \longrightarrow \oplus_{X_1} R^\infty  \stackrel{f_1}\longrightarrow R^\infty\stackrel{f_0} \longrightarrow R^\infty / \mathfrak m R^\infty \longrightarrow 0,
\]
where the \(X_i\) are sets non necessarily finite. Since \(R \subseteq R^\infty\), the above exact sequence is indeed \(R\)-linear.
Let $\Syz_j$ denote the $j$-th syzygy in this resolution of $R^\infty / \mathfrak m R^\infty$. Then the exact sequence
\[
0\longrightarrow \oplus_{X_t} R^\infty  
\longrightarrow \oplus_{X_{t-1}} R^\infty 
\longrightarrow \Syz_{\ell-1}
\longrightarrow 0
\]
induces a long exact sequence
\[
\operatorname{Ext}^{i}_{R}
(k,\oplus_{X_{t-1}} R^\infty)
\longrightarrow
\operatorname{Ext}^{i}_{R}(k,\Syz_{\ell-1})  \longrightarrow
\operatorname{Ext}^{i+1}_{R}
(k,\oplus_{X_t} R^\infty).
\]
Because $F_{t-1}:=\oplus_{X_{t-1}} R^\infty $ is free, and recall that direct sum is particularly a directed limit, then we have
\[
\operatorname{Ext}^{i}_{R}(k,F_{t-1})=\vil 	\operatorname{Ext}^{i}_{R}(k,R^\infty) =0
\qquad (\ast),
\] because 
 $R$  is coherent, $i>0$ and $k$
is  finitely presented (see {Fact} \ref{comuf}).
Hence the vanishing of $(\ast)$
implies
\[
\operatorname{Ext}^{i+1}_{R}
(k,\oplus_{X_t} R^\infty)=\vil	\operatorname{Ext}^{i+1}_{R}(k,R^\infty) =0,
\]here we need to assume that $i>\dim(R)$, see \cite[Theorem 1.2]{IYENGAR}.
From this,
\(
\operatorname{Ext}^{i}_{R}(k,\Omega_{\ell-1})=0 .
\)
Thus, for every $j$,
	there is a short exact sequence
	\[
	0\longrightarrow \Syz_{j}
	\longrightarrow F_{j-1}
	\longrightarrow \Syz_{j-1}
	\longrightarrow 0 .
	\]
	Hence by induction the vanishing of
	\(
	\operatorname{Ext}^{i}_{R}(k,\Syz_j)
	=0\)
	implies
	\(
	\operatorname{Ext}^{i+1}_{R}(k,\Syz_{j-1})=0 ,
	\)
	 in particular
	\(
	\operatorname{Ext}^{i}_{R}(k,k)=0 .
	\)
By Roberts' theorem (see \cite[Theorem 2]{Ro2}), we know $\mu^n(\mathfrak m,k)$ is nonzero only in the range
$
\bigl[\operatorname{depth}_R(k),\,\operatorname{id}_R(k)\bigr].	$ Therefore,
	\(
	\operatorname{id}_{R}(k)<\infty .
	\) Hence $R$ is a regular local ring.
\end{proof}

The following drops the isolated singularity from \cite[Theorem 3.15(4)]{Ryo} when we work with minimal perfect closure: 
\begin{corollary}\label{63}
Suppose 	\(
	\operatorname{Ext}^{+}_{R}(R^\infty,R^\infty)=0\).
Then 	$R$ is regular.
\end{corollary}


\begin{proof}
	Combine {Proposition} \ref{47} along with the previous theorem.
\end{proof}
Also, a similar argument  gives:

 \begin{corollary}(Aberbach-Li)
	Suppose 	\(
	\operatorname{Tor}^{R}_{i}(k,R^\infty)=0\)   $\exists i>0$.
	Then 	$R$ is regular.
\end{corollary}

The following is an easy approach to  Bhatt-Iyengar-Ma, see \cite{bim}. 
\begin{corollary}\label{63a}
	Suppose 	\(
	\operatorname{Tor}_{+}^{R}(R^\infty,R^\infty)=0\).
	Then 	$R$ is regular.
\end{corollary}


\begin{proof}
Against $\Ext$, tor modules commutes with localization. Then by an easy
induction on Krull dimension, we may assume in addition that $R$ is of isolated singularity, then an argument  similar to \cite[Theorem 3.15(4)]{Ryo} implies	$R$ is regular.
\end{proof}
It is worth noting that many of the results--Corollary \ref{63} among them--extend to 
$R^+$ as well. We do not pursue these extensions here, however, since our present focus is exclusively on 
$R^\infty$. Since our results are in prime characteristic, it is natural to ask (see \cite{MPa} and \cite{Ryo} for some progress):
\begin{question}Let $R$ be of characteristic zero.  When is $\operatorname{Gid}_R(R^+) < \infty$?
\end{question}
\begin{acknowledgement}
I thank everyone who helped me in carrying out this work. In particular, I thank Ryo Ishizuka for many valuable comments. I learned a lot about {Conjecture} \ref{61} with my coauthors Ishizuka \cite{Ryo}, Mahdavi \cite{Elham}, Patankar \cite{MPa} along with different projects. 
\end{acknowledgement}
\maketitle


	\begin{thebibliography}{99}

\bibitem{Ab}
I. M. Aberbach, J. Li, Asymptotic vanishing conditions which force regularity in local rings of prime
characteristic, Math. Res. Lett. {\bf 15}(4)  (2008), 815-820.

\bibitem{A}
M. Asgharzadeh, \emph{Homological properties of the perfect and absolute integral closure of Noetherian domains}, Math. Annalen {\bf 348} (2010), 237--263.
\bibitem{Elham}
M. Asgharzadeh, E. Mahdavi, \emph{Non-Noetherian Bass and Betti numbers},   arXiv:2606.20391.

\bibitem{Ryo}
M. Asgharzadeh, R. Ishizuka, 
\emph{Homological detection by perfectoid algebras},  	arXiv:2607.17801.


\bibitem{moh} M. Asgharzadeh,   \emph{Homological aspects of perfect algebras}, arXiv:1005.1462.

 
	\bibitem{moh5} M.
	Asgharzadeh,     \emph{Notes on modules of finite injective dimension}, arXiv:2301.01105.

		
\bibitem{MPa} M. Asgharzadeh,  S. Patankar, {\it Remarks on some homological problems regarding infinite integral extensions}, arXiv:2608.09473.

\bibitem{BS1}
B. Bhatt, and P. Scholze,  \emph{Projectivity of the Witt vector affine Grassmannian}, Invent. Math. {\bf{209}} (2017),  329--423.
	
\bibitem{bim}
B. Bhatt, S. B. Iyengar, and L. Ma, \emph{Regular Rings and Perfect(oid) Algebras}, Communications in Algebra, \textbf{47}(6) (2019), 2367--2383.

\bibitem{topo}	B.	Bhatt, M. Blickle, G. Lyubeznik, A. K. Singh, W.  Wenliang,
		\emph{Applications of perverse sheaves in commutative algebra},
		J. Reine Angew. Math. {\bf{825}} (2025), 81–138.
	
		
\bibitem{BH} W. Bruns and J. Herzog, \emph{Cohen-Macaulay rings}, Cambridge Studies in Advanced Mathematics, {\bf 39}, Cambridge University Press, Cambridge, 1993.

\bibitem{IYENGAR} L. W. Christensen, S.B. Iyengar, T. Marley,  
		\emph{Rigidity of Ext and Tor with coefficients in residue fields of a commutative Noetherian ring},
		Proc. Edinb. Math. Soc., II. Ser. {\bf 62}, No. 2, 305-321 (2019).


			\bibitem{Totally}
L.~W.~Christensen, 
 H. B. Foxby, and 
H. Holm, 
\emph{Beyond totally reflexive modules and back},  {Commutative Algebra}, ({2011}),  {{101}--{143}}.

		
\bibitem{G1}{S. Glaz}, {\it Commutative coherent rings},
Springer LNM, {\bf1371}, 1989.
		
	\bibitem{her} J. Herzog, \emph{Ringe der Charakteristik p und Frobenius funktoren}, Math. Z. 	{\bf140} (1974), 67-78.


		
		\bibitem{Mat}
	H. Matsumura, \emph{Commutative ring theory}, Cambridge studies in advanced math, \textbf{8}, (1986).
		
		\bibitem{Ro2} P. Roberts, \emph{Two applications of dualizing complexes over local rings}, Annales Scientifiques de l'École Normale Supérieure {\bf 9}(1), (1976) 103-106.

		
		
			
			\bibitem{Rot}J. J. Rotman,
		\emph{An introduction to homological algebra}, 
		Universitext
		Springer, New York, 2009, xiv+709 pp.

	\bibitem{sim}A.M.
	Simon, \emph{Some homological properties of complete modules}, Math. Proc. Camb. Phil. Soc.
	{\bf108}, 231–246 (1990).
	
			\bibitem{MP} L. Ma and T. Polstra, 	\emph{F-Singularities: A aommutative algebra approach},   available at https://www.math.purdue.edu/~lma/FSingularities.pdf.
	
	\bibitem{sa}S. Nasseh,	S. Sather-Wagstaff,
\emph{Cohen factorizations: weak functoriality and applications},   Journal of Pure and Applied Algebra, {\bf219} (2015), no. 3, 622-645.
	
	\bibitem{O}
	T. Ogoma, \emph{Noetherian property of inductive limits of Noetherian local rings},
	Proc. Japan Acad., Ser. A \textbf{{67}}(3), 68--69 (1991).
	
	
	\bibitem{PS}
	C. Peskine and L. Szpiro, \emph{Dimension projective finie et cohomologie locale},
	Publ. Math. IHES. 
	{\bf42} 
	(1973), 47-119
	
	
	
	\bibitem{Shi10}
	K. Shimomoto,
	\emph{$F$-coherent rings with applications to tight closure theory},
	J. Algebra {\bf338}, 24--34 (2011).
	\bibitem{Al}
	A. Wheeler, W. Zhang,
	\emph{Remarks on effective uniform Briancon-Skoda}, arXiv:2510.04004


	 	
	\end{thebibliography}
	
\end{document}	.	\bibitem{h} M. Hochster, 	\emph{Canonical elements in local cohomology modules and the direct summand conjecture}, J. Algebra {\bf84} (1983), no. 2, 503–553.		\bibitem{sc3} H. Schoutens, \emph{Projective dimension and the singular locus}, Comm. Algebra {\bf 31} (2003), 217–239.	
\bibitem{IYENGAR} L. W. Christensen, S.B. Iyengar, T. Marley,  
\emph{Rigidity of Ext and Tor with coefficients in residue fields of a commutative Noetherian ring},
Proc. Edinb. Math. Soc., II. Ser. {\bf 62}, No. 2, 305-321 (2019).
